\documentclass[11pt]{article}

\usepackage[margin=1in]{geometry}
\usepackage{amsmath,amssymb,amsthm,mathtools,mathrsfs,bm}
\usepackage{enumitem}
\usepackage{microtype}
\usepackage{hyperref}

\allowdisplaybreaks

\newtheorem{theorem}{Theorem}[section]
\newtheorem{proposition}[theorem]{Proposition}
\newtheorem{lemma}[theorem]{Lemma}
\newtheorem{corollary}[theorem]{Corollary}
\theoremstyle{definition}
\newtheorem{definition}[theorem]{Definition}
\theoremstyle{remark}
\newtheorem{remark}[theorem]{Remark}

\newcommand{\R}{\mathbb R}
\newcommand{\C}{\mathbb C}
\newcommand{\N}{\mathbb N}
\newcommand{\E}{\mathbb E}
\newcommand{\mcA}{\mathcal A}
\newcommand{\mcF}{\mathcal F}
\newcommand{\mcP}{\mathcal P}
\newcommand{\mcQ}{\mathcal Q}
\newcommand{\mcH}{\mathcal H}
\newcommand{\one}{\mathbf 1}
\newcommand{\dd}{\,\mathrm d}
\newcommand{\tr}{\operatorname{tr}}
\newcommand{\supp}{\operatorname{supp}}
\newcommand{\Ent}{\operatorname{Ent}}
\newcommand{\Var}{\operatorname{Var}}
\newcommand{\spec}{\operatorname{spec}}
\newcommand{\Law}{\operatorname{Law}}

\title{Thermodynamic Euclidean Construction of the Emergent $SU(3)$ Gauge Theory}
\author{}
\date{}

\begin{document}
\maketitle

\begin{abstract}
This paper implements the thermodynamic-first construction of the Euclidean
$SU(3)$ gauge theory used later in the series.  The Euclidean theory is defined
directly from the stationary mean-field equilibrium and the associated
stationary Markov path law furnished by the convergence paper, together with
the continuum-indexed gauge sector identified by the continuum-gauge paper.  We
package the exact Euclidean probability space, the local bounded observable
algebras on time windows, the continuum-indexed Schwinger functionals, the
thermodynamic identification of the gauge sector, the exact hyperplane Markov
factorization across the time-zero slice, and the transfer identity that
rewrites mixed Euclidean-time correlations as $L^2(\pi_\infty)$ semigroup
pairings.
\end{abstract}

\tableofcontents

\section{Introduction}

The construction uses three layers of structure:
\begin{enumerate}[leftmargin=1.5em]
\item the stationary thermodynamic particle theory of
\cite{convergence}, with unique mean-field equilibrium \(\pi_\infty\), its
stationary Markov path law, and the corresponding spectral-gap and entropy
information;
\item the deterministic flat regulator theory of \cite{continuumqft}, which
builds the finite Wilson sector from that thermodynamic substrate and identifies
its localized and global continuum Yang--Mills limits; and
\item the thermodynamic Euclidean theory constructed here from
that stationary substrate and continuum-indexed gauge sector.
\end{enumerate}

The Euclidean quantum theory is the stationary thermodynamic theory itself,
read through the continuum-indexed gauge sector proved in
\cite{continuumqft}.

The paper provides the realization-specific Euclidean machinery needed for the
axiomatic verification and reconstruction steps:
\begin{enumerate}[leftmargin=1.5em]
\item the stationary Euclidean probability space;
\item the local bounded observable algebras on time windows;
\item the continuum-indexed Schwinger functionals;
\item the constructive identification of the gauge sector in thermodynamic,
continuum, and deterministic finite-regulator presentations;
\item the hyperplane factorization and boundary-transfer formulas.
\end{enumerate}
These results are used directly in the verification of OS0--OS4 and in the
Osterwalder--Schrader reconstruction.

\begin{theorem}[Main thermodynamic Euclidean construction]
\label{thm:main}
The companion papers canonically determine a Euclidean \(SU(3)\) gauge-theory
data package
\[
\bigl(\Omega,\mcF,\mathbb P,\{\tau_t\}_{t\in\R},\Theta,\mcA_{\mathrm{loc}},
\omega_E,\mcA_{\mathrm{YM}}\bigr)
\]
with the following properties.
\begin{enumerate}[label=\textnormal{(\roman*)},leftmargin=1.5em]
\item \((\Omega,\mcF,\mathbb P)\) is the stationary Euclidean probability space
of the mean-field Markov theory; its one-time law is \(\pi_\infty\), its
time-translation group is \(\tau_t\), and \(\Theta\) is time reflection across
the hyperplane \(t=0\).
\item For every bounded time window \(I\subset\R\), there is a local bounded
observable algebra \(\mcA(I)\subset L^\infty(\Omega,\mcF_I,\mathbb P)\).  The
local Euclidean algebra is
\[
\mcA_{\mathrm{loc}}:=\bigcup_{I\Subset\R}\mcA(I).
\]
\item The continuum-indexed Schwinger functionals are the equilibrium moments
\[
\omega_E(F_1,\dots,F_n;t_1,\dots,t_n)
:=
\E_{\mathbb P}\!\bigl[(\tau_{t_1}F_1)\cdots(\tau_{t_n}F_n)\bigr],
\]
well defined for all \(n\ge1\), all \(F_j\in\mcA_{\mathrm{loc}}\), and all
times \(t_1,\dots,t_n\in\R\).
\item The distinguished gauge sector \(\mcA_{\mathrm{YM}}\subset\mcA_{\mathrm{loc}}\)
is indexed by continuum support data and contains the localized Wilson sector,
whose localized continuum Yang--Mills presentation is
\(S_{\mathrm{YM},\chi}^{\mathrm{cont}}[F^\infty]\) for
\(\chi\in C_c(\R^4)\), together with the path-space pullbacks of the one-time
and adjacent-time thermodynamic gauge generators identified in
\cite{continuumqft}.  The auxiliary finite-product plaquette functionals of
\cite{continuumqft} remain canonically attached through finite tensor powers of
the same stationary path law.  In that sense, the deterministic finite Wilson
sector, the continuum Yang--Mills functionals, and the thermodynamic
observables are compatible presentations of a common constructive gauge sector.
\item Across the reflection hyperplane \(t=0\), the positive-time and
negative-time sigma-algebras are conditionally independent given the time-zero
slice, and mixed past--future Euclidean correlations reduce to the stationary
mean-field semigroup on \(L^2(\pi_\infty)\) by an exact boundary transfer
formula.  The reflected positive-time case is a special instance of the same
identity.
\end{enumerate}
\end{theorem}

\section{Imported theorem package}

\begin{proposition}[Imported thermodynamic Yang--Mills package]
\label{prop:imports}
The companion papers establish the following facts.
\begin{enumerate}[label=\textnormal{(\alph*)},leftmargin=1.5em]
\item The convergence paper \cite{convergence} proves the existence and
uniqueness of the stationary mean-field law
\[
\pi_\infty\in\mcP(\R^8),
\]
the well posedness of the mean-field Markov dynamics for initial law
\(\pi_\infty\), the corresponding stationary Markov semigroup \((T_t)_{t\ge0}\)
on \(L^2(\pi_\infty)\), and a strict logarithmic Sobolev inequality for
\(\pi_\infty\).
\item The same convergence paper, together with the path-space bridge developed
in \cite{continuumqft}, yields the finite-horizon stationary path laws
\[
\Pi_{[0,S]}^\infty\in\mcP(C([0,S];\R^8)),
\qquad S>0,
\]
consistent under restriction and stationary under time shifts.
\item The continuum-gauge paper \cite{continuumqft} constructs the deterministic
flat regulator, its Wilson sector, the localized and global continuum
Yang--Mills functionals
\[
S_{\mathrm{YM},\chi}^{\mathrm{cont}}[F^\infty],
\qquad
S_{\mathrm{YM}}^{\mathrm{cont}}[F^\infty],
\]
and the associated thermodynamic link and plaquette observables on
\(\pi_\infty\) and its product laws.
\item The same continuum-gauge paper places those objects on the Euclidean
spacetime \(\R^4\): the deterministic flat regulator is embedded in
\(\R^4\), every localized continuum Yang--Mills functional is indexed by a test
function \(\chi\in C_c(\R^4)\), and the imported thermodynamic one-time and
adjacent-time gauge generators are furnished together with the corresponding
continuum support data on \(\R^4\).
\end{enumerate}
\end{proposition}

\begin{proof}
Item \textnormal{(a)} is the well-posedness, stationarity, and mean-field
entropy package proved in \cite{convergence}.  Item \textnormal{(b)} is the
finite-horizon stationary path-law package and stationary path-space
propagation-of-chaos framework recorded in \cite{continuumqft}.  Item
\textnormal{(c)} is the deterministic flat-regulator and Wilson-to-Yang--Mills
continuum-limit package of \cite{continuumqft}.  Item \textnormal{(d)} is the
support-labelled Euclidean spacetime package exported there: the regulator is
realized in \(\R^4\), the localized Yang--Mills functionals are indexed by
compact spacetime test functions, and the imported thermodynamic gauge
observables come with the same continuum support labels.
\end{proof}

\begin{remark}[Euclidean data package]
This section packages the imported results into the Euclidean thermodynamic data
package used in the axiomatic and reconstruction arguments.
\end{remark}

\section{Thermodynamic Euclidean probability space}

\begin{proposition}[Canonical two-sided stationary path law]
\label{prop:two-sided-extension}
For every bounded interval \([a,b]\subset\R\), set \(S:=b-a\) and define
\[
\sigma_a:C([0,S];\R^8)\to C([a,b];\R^8),
\qquad
(\sigma_a\gamma)(u):=\gamma(u-a).
\]
Let
\[
\Pi_{[a,b]}^\infty:=(\sigma_a)_\#\Pi_{[0,S]}^\infty.
\]
Then the family \(\{\Pi_{[a,b]}^\infty\}_{a<b}\) is projectively consistent under
restriction to smaller intervals.  Consequently, there exists a unique Borel
probability measure \(\mathbb P\) on \(\Omega:=C(\R;\R^8)\) such that for every
bounded interval \([a,b]\),
\[
(r_{[a,b]})_\#\mathbb P=\Pi_{[a,b]}^\infty,
\qquad
r_{[a,b]}(\omega):=\omega|_{[a,b]}.
\]
For \(t\in\R\), let
\[
(\tau_t\omega)(s):=\omega(s+t),
\qquad \omega\in\Omega,\ s\in\R.
\]
Moreover,
\[
\mathbb P\circ\tau_t^{-1}=\mathbb P
\qquad\forall\,t\in\R,
\]
and the one-time marginal of \(\mathbb P\) is \(\pi_\infty\).
\end{proposition}

\begin{proof}
By Proposition~\ref{prop:imports}\textnormal{(b)}, the finite-horizon laws
\(\Pi_{[0,S]}^\infty\) are consistent under restriction.  If
\([a',b']\subset[a,b]\), then
\[
(r_{[a',b']})_\#\Pi_{[a,b]}^\infty=\Pi_{[a',b']}^\infty.
\]
Thus \(\{\Pi_{[a,b]}^\infty\}_{a<b}\) forms a projectively consistent family on
the inverse system \(\{C([-m,m];\R^8)\}_{m\ge1}\). Since
\[
\Omega=C(\R;\R^8)\cong\varprojlim_{m\to\infty} C([-m,m];\R^8)
\]
as a Polish inverse limit, the standard projective-limit theorem yields a
unique Borel probability measure \(\mathbb P\) with the stated restrictions;
see, for example, \cite[Chapter~V]{parthasarathy1967}.  Equivalently, one may
apply the Kolmogorov extension theorem to the consistent compact-interval
cylinder laws and use the continuity built into the state spaces
\(C([-m,m];\R^8)\).
The one-time marginal is \(\pi_\infty\) because
\((\operatorname{ev}_0)_\#\Pi_{[0,S]}^\infty=\pi_\infty\), where
\(\operatorname{ev}_0(\gamma):=\gamma(0)\), by stationarity of the
imported finite-horizon laws.  Finally, for every bounded interval \([a,b]\),
\[
(r_{[a,b]})_\#(\mathbb P\circ\tau_t^{-1})
=
(\tau_t^{[a,b]})_\#\Pi_{[a+t,b+t]}^\infty,
\]
where \(\tau_t^{[a,b]}:C([a+t,b+t];\R^8)\to C([a,b];\R^8)\) is the restriction
shift \((\tau_t^{[a,b]}\eta)(u):=\eta(u+t)\). By the shift-stationarity imported
in Proposition~\ref{prop:imports}\textnormal{(b)},
\[
(\tau_t^{[a,b]})_\#\Pi_{[a+t,b+t]}^\infty=\Pi_{[a,b]}^\infty.
\]
Since bounded-interval restrictions determine Borel measures on \(\Omega\),
\(\mathbb P\) is stationary.
\end{proof}

\begin{definition}[Stationary Euclidean probability space]
\label{def:euclidean-space}
Let \(\Omega:=C(\R;\R^8)\) be the canonical two-sided path space, equipped with
the coordinate process
\[
Y_t(\omega):=\omega(t),
\qquad t\in\R.
\]
Let \(\mcF\) be its Borel sigma-algebra and let \(\mathbb P\) be the canonical
two-sided stationary path law from
Proposition~\ref{prop:two-sided-extension}.  For \(t\in\R\), define the
time-translation map
\[
(\tau_t\omega)(s):=\omega(s+t),
\qquad s\in\R,
\]
and let \(\Theta\) denote time reflection across \(t=0\),
\[
(\Theta\omega)(s):=\omega(-s).
\]
For each Borel set \(I\subset\R\), set
\[
\mcF_I:=\sigma\bigl(Y_s:s\in I\bigr).
\]
In particular,
\[
\mcF_+:=\mcF_{(0,\infty)},
\qquad
\mcF_-:=\mcF_{(-\infty,0)},
\qquad
\mcF_0:=\mcF_{\{0\}}.
\]
When \(\tau_t\) and \(\Theta\) act on observables, we use the pullback
convention
\[
(\tau_tF)(\omega):=F(\tau_t\omega),
\qquad
(\Theta F)(\omega):=F(\Theta\omega).
\]
\end{definition}

\begin{remark}
The measure \(\mathbb P\) is stationary by construction:
\[
\mathbb P\circ \tau_t^{-1}=\mathbb P
\qquad
\forall\,t\in\R.
\]
No axiom verification is being made at this stage.  Definition
\ref{def:euclidean-space} merely fixes the Euclidean probability space on which
all later Schwinger functionals are evaluated.
\end{remark}

\begin{definition}[Local bounded observable algebras on time windows]
\label{def:local-algebra}
For every bounded interval \(I\Subset\R\), define
\[
\mcA(I):=L^\infty(\Omega,\mcF_I,\mathbb P).
\]
Set
\[
\mcA_{\mathrm{loc}}:=\bigcup_{I\Subset\R}\mcA(I).
\]
An observable in \(\mcA_{\mathrm{loc}}\) is called \emph{local} if it belongs
to \(\mcA(I)\) for some bounded interval \(I\).
\end{definition}

\begin{proposition}[Basic algebraic properties of the local bounded observable algebras]
\label{prop:local-algebra}
The local bounded observable algebras satisfy the following.
\begin{enumerate}[label=\textnormal{(\roman*)},leftmargin=1.5em]
\item For every bounded interval \(I\), \(\mcA(I)\) is a unital \( * \)-algebra.
\item If \(I_1\subset I_2\), then \(\mcA(I_1)\subset\mcA(I_2)\).
\item For every \(t\in\R\), \(\tau_t\) maps \(\mcA(I)\) onto
\(\mcA(I+t)\), and \(\Theta\) maps \(\mcA(I)\) onto \(\mcA(-I)\).
\end{enumerate}
\end{proposition}

\begin{proof}
Each \(\mcA(I)\) is an \(L^\infty\)-algebra over the sigma-algebra \(\mcF_I\),
hence is a unital \( * \)-algebra.  If \(I_1\subset I_2\), then
\(\mcF_{I_1}\subset\mcF_{I_2}\), so \(L^\infty(\mcF_{I_1})\subset
L^\infty(\mcF_{I_2})\).  If \(F\) is \(\mcF_I\)-measurable, then
\((\tau_tF)(\omega)=F(\omega(\cdot+t))\) depends only on the coordinates
\(\omega(s)\) with \(s\in I+t\), hence \(\tau_tF\in\mcA(I+t)\).  The statement
for \(\Theta\) is identical.
\end{proof}

\begin{definition}[Stationary Euclidean state]
\label{def:euclidean-state}
The stationary Euclidean state on \(\mcA_{\mathrm{loc}}\) is the equilibrium
expectation functional
\[
\omega_E(F):=\E_{\mathbb P}[F],
\qquad F\in\mcA_{\mathrm{loc}}.
\]
\end{definition}

\begin{definition}[Schwinger functionals]
\label{def:schwinger}
For \(n\ge1\), local bounded observables \(F_1,\dots,F_n\in\mcA_{\mathrm{loc}}\),
and times \(t_1,\dots,t_n\in\R\), define the \(n\)-point Schwinger functional
by
\[
\omega_E(F_1,\dots,F_n;t_1,\dots,t_n)
:=
\E_{\mathbb P}\!\bigl[(\tau_{t_1}F_1)\cdots(\tau_{t_n}F_n)\bigr].
\]
\end{definition}

\begin{proposition}[Well-definedness and stationarity of the Schwinger functionals]
\label{prop:schwinger-basic}
The Schwinger functionals of Definition~\ref{def:schwinger} are well defined
for all local bounded observables.  They are multilinear and satisfy the bound
\[
\left|\omega_E(F_1,\dots,F_n;t_1,\dots,t_n)\right|
\le
\prod_{j=1}^n \|F_j\|_{L^\infty(\mathbb P)}.
\]
Moreover, for every \(s\in\R\),
\[
\omega_E(F_1,\dots,F_n;t_1+s,\dots,t_n+s)
=
\omega_E(F_1,\dots,F_n;t_1,\dots,t_n).
\]
\end{proposition}

\begin{proof}
If \(F_j\in\mcA(I_j)\), then \(\tau_{t_j}F_j\in
L^\infty(\Omega,\mcF_{I_j+t_j},\mathbb P)\), hence the product is bounded and
measurable.  The norm bound is immediate from
\[
\bigl|(\tau_{t_1}F_1)\cdots(\tau_{t_n}F_n)\bigr|
\le
\prod_{j=1}^n \|F_j\|_{L^\infty(\mathbb P)}.
\]
Multilinearity is immediate from linearity of expectation.  Stationarity of
\(\mathbb P\) gives
\[
\E_{\mathbb P}\!\bigl[(\tau_{t_1+s}F_1)\cdots(\tau_{t_n+s}F_n)\bigr]
=
\E_{\mathbb P}\!\bigl[(\tau_{t_1}F_1)\cdots(\tau_{t_n}F_n)\bigr].
\]
\end{proof}

\begin{definition}[Finite-time evaluation maps]
\label{def:evaluation-maps}
For \(k\ge1\) and times \(t_1,\dots,t_k\in\R\), define
\[
\operatorname{ev}_{t_1,\dots,t_k}:\Omega\to(\R^8)^k,
\qquad
\operatorname{ev}_{t_1,\dots,t_k}(\omega)
:=
\bigl(Y_{t_1}(\omega),\dots,Y_{t_k}(\omega)\bigr).
\]
If \(h:(\R^8)^k\to\C\) is bounded and Borel, write
\[
\operatorname{ev}_{t_1,\dots,t_k}^*h:=h\circ\operatorname{ev}_{t_1,\dots,t_k}.
\]
\end{definition}

\begin{proposition}[Locality of finite-time pullbacks]
\label{prop:evaluation-maps}
For every bounded Borel \(h:(\R^8)^k\to\C\) and every times
\(t_1,\dots,t_k\in\R\),
\[
\operatorname{ev}_{t_1,\dots,t_k}^*h
\in
\mcA\bigl([\min_j t_j,\max_j t_j]\bigr).
\]
In particular, every such pullback belongs to \(\mcA_{\mathrm{loc}}\).
\end{proposition}

\begin{proof}
Each coordinate map \(Y_{t_j}\) is Borel, hence
\(\operatorname{ev}_{t_1,\dots,t_k}\) is Borel.  The pullback depends only on
the finitely many coordinates \(Y_{t_1},\dots,Y_{t_k}\), so it is measurable
with respect to
\[
\sigma(Y_{t_1},\dots,Y_{t_k})
\subset
\mcF_{[\min_j t_j,\max_j t_j]}.
\]
Boundedness of \(h\) gives boundedness of the pullback.
\end{proof}

\section{Constructive identification of the gauge sector}

\begin{definition}[Ambient and physical Euclidean observable algebras]
\label{def:ambient-physical}
The algebra
\[
\mcA_{\mathrm{loc}}=\bigcup_{I\Subset\R}L^\infty(\Omega,\mcF_I,\mathbb P)
\]
is the ambient local measurable algebra on the thermodynamic path space.
The physical Euclidean observable algebra of the Yang--Mills theory realized in
this paper is the distinguished gauge sector
\[
\mcA_{\mathrm{phys}}:=\mcA_{\mathrm{YM}}\subset\mcA_{\mathrm{loc}}.
\]
Thus \(\mcA_{\mathrm{loc}}\) records all bounded path-space local observables,
whereas \(\mcA_{\mathrm{phys}}\) records the physical gauge-invariant
subalgebra selected by the constructive series.
\end{definition}

\begin{definition}[Distinguished thermodynamic gauge sector]
\label{def:gauge-sector}
Let \(\mcA_{\mathrm{YM}}\subset\mcA_{\mathrm{loc}}\) be the smallest unital
\(*\)-subalgebra containing the following continuum-indexed generators:
\begin{enumerate}[label=\textnormal{(\roman*)},leftmargin=1.5em]
\item for every \(\chi\in C_c(\R^4)\), the constant observable
\[
\omega\longmapsto S_{\mathrm{YM},\chi}^{\mathrm{cont}}[F^\infty];
\]
\item the constant observable
\[
\omega\longmapsto S_{\mathrm{YM}}^{\mathrm{cont}}[F^\infty];
\]
\item every pullback
\[
\operatorname{ev}_{t_1,\dots,t_k}^*h
\]
for which \(h\) is a bounded Borel representative of one of the one-time or
adjacent-time thermodynamic gauge generators furnished by
\cite{continuumqft}, and \((t_1,\dots,t_k)\) is the corresponding finite-time
index set.
\end{enumerate}
The additional finite-product plaquette functionals imported from
\cite{continuumqft} are canonically attached to the same stationary substrate
through finite tensor powers \(\mathbb P^{\otimes m}\) of the path law and are
used here only at the level of Euclidean correlation functionals; they are not
introduced as extra pointwise generators of the one-copy algebra
\(\mcA_{\mathrm{YM}}\).
\end{definition}

\begin{proposition}[Physical generating family inside the ambient algebra]
\label{prop:physical-generators}
Every realization-specific Yang--Mills observable imported from
Proposition~\ref{prop:imports}\textnormal{(c)--(d)} and used to define the
Euclidean gauge theory belongs to the distinguished physical algebra
\(\mcA_{\mathrm{YM}}\subset\mcA_{\mathrm{loc}}\).  More precisely:
\begin{enumerate}[label=\textnormal{(\roman*)},leftmargin=1.5em]
\item for every \(\chi\in C_c(\R^4)\), the localized and global continuum
Yang--Mills functionals define bounded constant observables in
\(\mcA_{\mathrm{YM}}\);
\item every imported one-time or adjacent-time thermodynamic gauge generator
admits a bounded Borel representative whose path-space pullback belongs to
\(\mcA_{\mathrm{YM}}\subset\mcA_{\mathrm{loc}}\);
\item \(\mcA_{\mathrm{YM}}\) is the smallest unital \( * \)-subalgebra of
\(\mcA_{\mathrm{loc}}\) containing that realization-specific physical
generating family.
\end{enumerate}
\end{proposition}

\begin{proof}
Item \textnormal{(i)} is part of Definition~\ref{def:gauge-sector}.  Item
\textnormal{(ii)} follows from Proposition~\ref{prop:evaluation-maps}, which
places the corresponding bounded Borel pullbacks in \(\mcA_{\mathrm{loc}}\),
and from Definition~\ref{def:gauge-sector}, which includes exactly those
pullbacks in \(\mcA_{\mathrm{YM}}\).  Item \textnormal{(iii)} is the minimality
statement built into Definition~\ref{def:gauge-sector}.
\end{proof}

\begin{theorem}[Thermodynamic-first identification of the Euclidean gauge theory]
\label{thm:gauge-identification}
The Euclidean gauge theory used in the series is the thermodynamic
theory defined by the stationary path law \(\mathbb P\) of
Definition~\ref{def:euclidean-space}, together with the distinguished
continuum-indexed gauge sector \(\mcA_{\mathrm{YM}}\) of
Definition~\ref{def:gauge-sector}.  More precisely:
\begin{enumerate}[label=\textnormal{(\roman*)},leftmargin=1.5em]
\item the foundational Euclidean state is the stationary equilibrium state
\(\omega_E\) of Definition~\ref{def:euclidean-state};
\item the gauge sector is read through the continuum Yang--Mills functional
\(S_{\mathrm{YM}}^{\mathrm{cont}}[F^\infty]\), the path-space pullbacks of the
imported one-time and adjacent-time thermodynamic gauge generators, and the
canonically attached finite-product plaquette functionals of
\cite{continuumqft};
\item the deterministic flat-regulator Wilson sector of \cite{continuumqft} is
a compatible finite presentation on the same emergent geometry, but it is not
used here as the defining approximation scheme for the continuum state.
\end{enumerate}
\end{theorem}

\begin{proof}
Item \textnormal{(i)} is just the thermodynamic-first definition of the
Euclidean theory adopted here: the state \(\omega_E\) is the
equilibrium expectation under the stationary path law \(\mathbb P\).
Item \textnormal{(ii)} is the continuum-gauge identification provided by
Proposition~\ref{prop:imports}\textnormal{(c)} together with
Proposition~\ref{prop:evaluation-maps}.  Item \textnormal{(iii)} is the finite
Wilson presentation from the same deterministic flat-regulator package.  These
statements are mutually compatible because they are all built on the same
deterministic flat regulator family selected by the thermodynamic particle
substrate.
\end{proof}

\begin{corollary}[Finite-regulator, continuum, and thermodynamic presentations of a common constructive gauge sector]
\label{cor:three-presentations}
The \(SU(3)\) gauge sector admits three compatible presentations:
\begin{enumerate}[label=\textnormal{(\roman*)},leftmargin=1.5em]
\item the exact finite-regulator Wilson theory on the deterministic flat
regulators;
\item the continuum Yang--Mills functional
\(S_{\mathrm{YM}}^{\mathrm{cont}}[F^\infty]\);
\item the thermodynamic one-time and adjacent-time gauge generators, together
with the canonically attached finite-product plaquette functionals, on the
stationary substrate determined by \(\pi_\infty\) and the path laws imported
from \cite{continuumqft}.
\end{enumerate}
The third presentation, together with its continuum Yang--Mills
identification, is taken as the foundational Euclidean theory.
\end{corollary}

\begin{proof}
The compatibility of the three presentations is the content of
Theorem~\ref{thm:gauge-identification}.
\end{proof}

\begin{theorem}[Exact identification of the physical Euclidean gauge sector]
\label{thm:exact-gauge-sector}
The physical Euclidean gauge sector is the
distinguished continuum-indexed algebra \(\mcA_{\mathrm{YM}}\) of
Definition~\ref{def:gauge-sector} together with the state
\(\omega_E\) of Definition~\ref{def:schwinger}.  Equivalently, within the
thermodynamic Euclidean realization, the following describe the same physical
gauge sector:
\begin{enumerate}[label=\textnormal{(\roman*)},leftmargin=1.5em]
\item the deterministic finite Wilson presentation on the flat regulators;
\item the continuum Yang--Mills functionals
\(S_{\mathrm{YM}}^{\mathrm{cont}}[F^\infty]\) and
\(S_{\mathrm{YM},\chi}^{\mathrm{cont}}[F^\infty]\);
\item the imported one-time and adjacent-time thermodynamic gauge generators,
together with the canonically attached finite-product plaquette functionals;
\item the Euclidean local gauge sector obtained by evaluating the corresponding
continuum-indexed observables on the stationary path law \(\mathbb P\).
\end{enumerate}
These descriptions are equivalent presentations of the same physical
Euclidean \(SU(3)\) gauge theory.
\end{theorem}

\begin{proof}
Theorem~\ref{thm:gauge-identification} identifies the Euclidean gauge theory
with the stationary path law \(\mathbb P\) together with the
distinguished gauge sector \(\mcA_{\mathrm{YM}}\).  Corollary~\ref{cor:three-presentations}
states that the deterministic finite Wilson theory, the continuum
Yang--Mills functional, and the thermodynamic observable presentation are
compatible presentations of that same gauge sector.  Since
Definition~\ref{def:gauge-sector} takes \(\mcA_{\mathrm{YM}}\) to be the
smallest unital \( * \)-subalgebra containing exactly those continuum-indexed
generators, the physical Euclidean gauge sector is precisely the one stated
above.
\end{proof}

\begin{proposition}[Gauge-sector closure]
\label{prop:gauge-sector-closure}
The distinguished gauge sector \(\mcA_{\mathrm{YM}}\) is a unital
\( * \)-subalgebra of \(\mcA_{\mathrm{loc}}\), is stable under Euclidean-time
translation and time reflection, and is indexed entirely by continuum support
data.
\end{proposition}

\begin{proof}
By Definition~\ref{def:gauge-sector}, \(\mcA_{\mathrm{YM}}\) is the smallest
unital \( * \)-subalgebra of \(\mcA_{\mathrm{loc}}\) containing the listed
continuum-indexed generators, so it is closed under products, linear
combinations, adjoints, and constants by construction.  Proposition~\ref{prop:local-algebra}\textnormal{(iii)}
shows that \(\tau_t\) and \(\Theta\) preserve locality on
\(\mcA_{\mathrm{loc}}\), and the generators of \(\mcA_{\mathrm{YM}}\) are
defined by continuum support data rather than regulator labels; therefore their
translated and reflected images remain in the same continuum-indexed gauge
sector.  Proposition~\ref{prop:continuum-indices} gives the final claim.
\end{proof}

\begin{proposition}[Wilson and thermodynamic generators determine the gauge sector]
\label{prop:wilson-generation}
The algebra \(\mcA_{\mathrm{YM}}\) is generated, in the algebraic sense of
Definition~\ref{def:gauge-sector}, by the localized Wilson observables,
the continuum Yang--Mills functionals, and the imported one-time and
adjacent-time thermodynamic gauge generators.  In particular, no additional
primitive gauge-sector generators are introduced beyond this Yang--Mills /
Wilson presentation.
\end{proposition}

\begin{proof}
This is how \(\mcA_{\mathrm{YM}}\) is defined in
Definition~\ref{def:gauge-sector}: it is the smallest unital \( * \)-subalgebra
containing those generators and no others.  Therefore every element of
\(\mcA_{\mathrm{YM}}\) is obtained from them by finite algebraic operations.
\end{proof}

\begin{definition}[Local physical gauge generators and global scalar sector]
\label{def:local-global-gauge}
Let \(\mcA_{\mathrm{YM}}^{\mathrm{loc}}\) denote the unital \( * \)-subalgebra
of \(\mcA_{\mathrm{YM}}\) generated by:
\begin{enumerate}[label=\textnormal{(\roman*)},leftmargin=1.5em]
\item the localized continuum Yang--Mills generators
\[
\omega\longmapsto S_{\mathrm{YM},\chi}^{\mathrm{cont}}[F^\infty],
\qquad \chi\in C_c(\R^4),
\]
\item the path-space pullbacks of the imported one-time and adjacent-time
thermodynamic gauge generators.
\end{enumerate}
Let
\[
Z_{\mathrm{YM}}
:=
S_{\mathrm{YM}}^{\mathrm{cont}}[F^\infty]\mathbf 1
\]
denote the global continuum Yang--Mills scalar, viewed as a constant element of
\(\mcA_{\mathrm{YM}}\).
\end{definition}

\begin{proposition}[Local-global decomposition of the physical gauge algebra]
\label{prop:local-global-decomposition}
The physical gauge algebra splits into its local gauge sector and the adjunction
of the global central Yang--Mills scalar:
\[
\mcA_{\mathrm{YM}}
=
\operatorname{alg}\bigl(\mcA_{\mathrm{YM}}^{\mathrm{loc}},Z_{\mathrm{YM}}\bigr).
\]
Moreover:
\begin{enumerate}[label=\textnormal{(\roman*)},leftmargin=1.5em]
\item every element of \(\mcA_{\mathrm{YM}}^{\mathrm{loc}}\) is generated by
continuum-indexed local gauge generators with nontrivial support data in
\(\R^4\);
\item \(Z_{\mathrm{YM}}\) is central in \(\mcA_{\mathrm{YM}}\);
\item the regional net \(O\mapsto\mcA_{\mathrm{YM}}(O)\) is generated entirely
from the local sector \(\mcA_{\mathrm{YM}}^{\mathrm{loc}}\), while
\(Z_{\mathrm{YM}}\) belongs to the global physical algebra \(\mcA_{\mathrm{YM}}\)
as a support-free scalar.
\end{enumerate}
\end{proposition}

\begin{proof}
Definition~\ref{def:gauge-sector} shows that \(\mcA_{\mathrm{YM}}\) is generated
by three families: the localized continuum Yang--Mills generators, the global
continuum Yang--Mills scalar, and the imported thermodynamic gauge pullbacks.
By Definition~\ref{def:local-global-gauge}, the first and third families
generate \(\mcA_{\mathrm{YM}}^{\mathrm{loc}}\), while the second is exactly
\(Z_{\mathrm{YM}}\).  This proves the displayed algebraic identity.

Item \textnormal{(i)} is immediate from the definition of
\(\mcA_{\mathrm{YM}}^{\mathrm{loc}}\).  Item \textnormal{(ii)} holds because
\(Z_{\mathrm{YM}}\) is a constant observable on \(\Omega\), hence commutes with
every multiplication operator in \(\mcA_{\mathrm{YM}}\subset L^\infty(\Omega,\mcF,\mathbb P)\).
For item \textnormal{(iii)}, Definition~\ref{def:regional-gauge-net} selects
generators by continuum support data contained in \(O\subset\R^4\).  The
support-carrying generators are precisely those in
\(\mcA_{\mathrm{YM}}^{\mathrm{loc}}\), whereas \(Z_{\mathrm{YM}}\) is a global
support-free scalar.
\end{proof}

\begin{proposition}[Continuum indexation of the gauge sector]
\label{prop:continuum-indices}
The distinguished gauge sector \(\mcA_{\mathrm{YM}}\) is indexed by continuum
support data rather than by regulator labels.  In particular:
\begin{enumerate}[label=\textnormal{(\roman*)},leftmargin=1.5em]
\item localized Wilson generators are indexed by test functions
\(\chi\in C_c(\R^4)\);
\item path-space gauge generators are indexed by the continuum support data and
finite-time arguments appearing in the imported one-time and adjacent-time
thermodynamic observables;
\item the auxiliary finite-product plaquette functionals are indexed by the same
continuum support data together with the corresponding finite tensor-power
level;
\item the Euclidean Schwinger functionals restricted to \(\mcA_{\mathrm{YM}}\)
therefore depend on continuum support data and smearing functions, not on the
regulator label \(N\).
\end{enumerate}
\end{proposition}

\begin{proof}
By Proposition~\ref{prop:imports}\textnormal{(c)}, the localized Wilson sector
is indexed by \(\chi\in C_c(\R^4)\), and the imported thermodynamic gauge
generators carry continuum support labels.  Proposition~\ref{prop:evaluation-maps}
turns their bounded Borel representatives into local observables on path space.
The finite-regulator label \(N\) appears only in the auxiliary approximation
statements that identify these observables with their regulator presentations.
Once the thermodynamic-first Euclidean state is fixed, the remaining labels are
exactly the continuum support data.
\end{proof}

\begin{definition}[Euclidean transport of support data]
\label{def:support-transport}
For \(a\in\R^4\), write
\[
\vartheta_a(x):=x+a,
\qquad x\in\R^4,
\]
and let \(\rho\) denote time reflection on \(\R^4\),
\[
\rho(x_0,x_1,x_2,x_3):=(-x_0,x_1,x_2,x_3).
\]
If \(\chi\in C_c(\R^4)\), define the transported test functions
\[
(\vartheta_a\chi)(x):=\chi(x-a),
\qquad
(\rho\chi)(x):=\chi(\rho x).
\]
For a continuum support datum \(s\) attached to one of the imported
thermodynamic gauge generators, write \(\vartheta_a s\) and \(\rho s\) for the
support data obtained by applying the same Euclidean maps to the underlying
spacetime support labels.
\end{definition}

\begin{proposition}[Euclidean support transport for gauge generators]
\label{prop:support-transport}
The continuum-indexed generators of \(\mcA_{\mathrm{YM}}\) are stable under the
Euclidean transport of their support data in the following sense.
\begin{enumerate}[label=\textnormal{(\roman*)},leftmargin=1.5em]
\item For every \(\chi\in C_c(\R^4)\) and every \(a\in\R^4\), the transported
localized Yang--Mills functionals
\[
S_{\mathrm{YM},\vartheta_a\chi}^{\mathrm{cont}}[F^\infty]
\qquad\text{and}\qquad
S_{\mathrm{YM},\rho\chi}^{\mathrm{cont}}[F^\infty]
\]
are again generators of \(\mcA_{\mathrm{YM}}\).
\item If
\[
\operatorname{ev}_{t_1,\dots,t_k}^*h
\]
is a generator coming from an imported one-time or adjacent-time thermodynamic
gauge observable with continuum support datum \(s\), then the generators with
transported support data \(\vartheta_a s\) and \(\rho s\) are again among the
imported continuum-indexed generators of \(\mcA_{\mathrm{YM}}\).
\item Consequently, if a gauge generator has support data contained in a bounded
region \(O\subset\R^4\), then the generators obtained by translating or
reflecting its support data have support contained in \(\vartheta_a(O)\) or
\(\rho(O)\), respectively.
\end{enumerate}
\end{proposition}

\begin{proof}
Item \textnormal{(i)} follows from
Proposition~\ref{prop:imports}\textnormal{(d)} together with
Definition~\ref{def:gauge-sector}: every compactly supported test function on
\(\R^4\) indexes a localized continuum Yang--Mills generator, and
\(\vartheta_a\chi,\rho\chi\in C_c(\R^4)\) whenever \(\chi\in C_c(\R^4)\).

For item \textnormal{(ii)}, Proposition~\ref{prop:imports}\textnormal{(d)}
states that the imported thermodynamic one-time and adjacent-time gauge
generators are furnished together with continuum support data on \(\R^4\).
Applying \(\vartheta_a\) or \(\rho\) to those support labels therefore produces
new continuum support data of the same type, hence generators of the same
imported family.  Definition~\ref{def:gauge-sector} includes every such
continuum-indexed imported generator.

Item \textnormal{(iii)} is immediate because Euclidean translations and time
reflection send subsets of \(\R^4\) to their translated or reflected images and
the transported generators remain in the same continuum-indexed family by
items~\textnormal{(i)} and \textnormal{(ii)}.
\end{proof}

\begin{definition}[Regional physical gauge net]
\label{def:regional-gauge-net}
For every bounded Euclidean region \(O\subset\R^4\), let
\[
\mcA_{\mathrm{YM}}(O)
\]
denote the unital \( * \)-subalgebra of \(\mcA_{\mathrm{YM}}\) generated by
those continuum-indexed gauge observables whose support data are contained in
\(O\).
\end{definition}

\begin{proposition}[Locality of the regional physical gauge net]
\label{prop:regional-locality}
For any bounded Euclidean regions \(O_1,O_2\subset\R^4\) and any
\[
F\in\mcA_{\mathrm{YM}}(O_1),\qquad G\in\mcA_{\mathrm{YM}}(O_2),
\]
one has
\[
FG=GF.
\]
In particular, the regional physical gauge net is local in the commutative
Euclidean sense.
\end{proposition}

\begin{proof}
By Definition~\ref{def:regional-gauge-net},
\(\mcA_{\mathrm{YM}}(O_j)\subset\mcA_{\mathrm{YM}}\subset\mcA_{\mathrm{loc}}
\subset L^\infty(\Omega,\mcF,\mathbb P)\) for \(j=1,2\).  Every element of
\(L^\infty(\Omega,\mcF,\mathbb P)\) acts by multiplication on the underlying
probability space, and multiplication of bounded measurable functions is
commutative.  Therefore \(FG=GF\).
\end{proof}

\begin{theorem}[Local net realization of the Yang--Mills gauge sector]
\label{thm:local-net-equivalence}
The assignment \(O\mapsto \mcA_{\mathrm{YM}}(O)\) is the physical Euclidean
local gauge net.  More precisely:
\begin{enumerate}[label=\textnormal{(\roman*)},leftmargin=1.5em]
\item if \(O_1\subset O_2\), then \(\mcA_{\mathrm{YM}}(O_1)\subset
\mcA_{\mathrm{YM}}(O_2)\);
\item \(\mcA_{\mathrm{YM}}(O)\subset\mcA_{\mathrm{YM}}\subset\mcA_{\mathrm{loc}}\)
for every bounded \(O\);
\item for any bounded regions \(O_1,O_2\subset\R^4\), every
\(F\in\mcA_{\mathrm{YM}}(O_1)\) commutes with every
\(G\in\mcA_{\mathrm{YM}}(O_2)\);
\item the net is stable under Euclidean translations and time reflection in the
sense that translated or reflected support data produce the corresponding
translated or reflected local algebra.
\end{enumerate}
Hence the physical Euclidean gauge theory is identified region by region by its
Yang--Mills local observable net, not only by a global action functional.
\end{theorem}

\begin{proof}
By Definition~\ref{def:regional-gauge-net}, \(\mcA_{\mathrm{YM}}(O)\) is
generated by those elements of \(\mcA_{\mathrm{YM}}\) whose continuum support
data lie in \(O\), so \(\mcA_{\mathrm{YM}}(O)\subset\mcA_{\mathrm{YM}}\subset
\mcA_{\mathrm{loc}}\).  If \(O_1\subset O_2\), every generator supported in
\(O_1\) is also supported in \(O_2\), which gives isotony.
Proposition~\ref{prop:regional-locality} gives commutativity of different
regional algebras.  Proposition~\ref{prop:gauge-sector-closure}
and Proposition~\ref{prop:continuum-indices} show that the sector is indexed by
continuum support data, while Proposition~\ref{prop:support-transport} gives
stability of the generator family under Euclidean transport of those support
labels.  Proposition~\ref{prop:local-global-decomposition} identifies the
support-carrying local sector from which the regional net is generated.
Therefore the resulting regional assignment is the local
Yang--Mills gauge net encoded by the realization.
\end{proof}

\begin{theorem}[No extra physical degrees of freedom beyond the Yang--Mills sector]
\label{thm:no-extra-sector}
Within the thermodynamic Euclidean realization, no
additional physical local observable sector survives beyond the distinguished
continuum Yang--Mills gauge sector.  Equivalently, every physical local
observable used in the Euclidean theory belongs to the gauge-sector net
\(O\mapsto\mcA_{\mathrm{YM}}(O)\) generated by the continuum Yang--Mills,
Wilson, and thermodynamic gauge observables.
\end{theorem}

\begin{proof}
Definition~\ref{def:ambient-physical} distinguishes the ambient measurable
algebra \(\mcA_{\mathrm{loc}}\) from the physical Euclidean algebra
\(\mcA_{\mathrm{phys}}=\mcA_{\mathrm{YM}}\).  Theorem~\ref{thm:gauge-identification}
fixes the Euclidean gauge theory of the series to be the stationary path law
\(\mathbb P\) together with that physical algebra.  Proposition~\ref{prop:physical-generators}
shows that every realization-specific Yang--Mills observable used to define the
physical Euclidean theory lies in \(\mcA_{\mathrm{YM}}\), and that
\(\mcA_{\mathrm{YM}}\) is the smallest unital \( * \)-subalgebra of
\(\mcA_{\mathrm{loc}}\) containing that physical generating family.
Proposition~\ref{prop:local-global-decomposition} identifies the local and
global parts of this physical algebra, and Theorem~\ref{thm:local-net-equivalence}
identifies the corresponding regional net.  Therefore no second independent
physical local observable sector survives beyond \(\mcA_{\mathrm{YM}}\): the
physical sector is exactly the Yang--Mills gauge sector already identified.
\end{proof}

\begin{corollary}[Regulator elimination at the level of the physical Euclidean theory]
\label{cor:regulator-elimination}
After passage to the thermodynamic Euclidean gauge sector, the resulting
Schwinger functionals and local gauge algebras depend only on the continuum
\(SU(3)\) curvature profile \(F^\infty\), the associated continuum support
data, and the stationary law \(\mathbb P\), and not on the auxiliary
finite-regulator label \(N\).
\end{corollary}

\begin{proof}
Proposition~\ref{prop:continuum-indices} shows that the gauge sector is indexed
by continuum support data rather than regulator labels.  Theorem~\ref{thm:exact-gauge-sector}
identifies the physical Euclidean gauge theory with that continuum-indexed
sector on the stationary path law.  Therefore the regulator label \(N\) enters
only through the auxiliary approximation results used upstream to identify the
continuum sector; it is not part of the final physical Euclidean theory.
\end{proof}

\section{Hyperplane factorization and boundary transfer}

\begin{proposition}[Exact hyperplane Markov factorization]
\label{prop:hyperplane-factorization}
In the stationary Euclidean probability space of
Definition~\ref{def:euclidean-space}, the positive-time and negative-time
sigma-algebras are conditionally independent given \(\mcF_0\).  Equivalently,
for bounded \(\mcF_+\)-measurable \(F\) and bounded \(\mcF_-\)-measurable \(G\),
\[
\E_{\mathbb P}[FG\mid\mcF_0]
=
\E_{\mathbb P}[F\mid\mcF_0]\,
\E_{\mathbb P}[G\mid\mcF_0].
\]
\end{proposition}

\begin{proof}
By Proposition~\ref{prop:imports}\textnormal{(a)--(b)}, \(\mathbb P\) is the
stationary path law of a Markov process with one-time marginal \(\pi_\infty\).
The Markov property implies that future and past are conditionally independent
given the present-time configuration \(Y_0\).  Since
\(\mcF_0=\sigma(Y_0)\), this is exactly the stated conditional-independence
identity.
\end{proof}

\begin{proposition}[Stability under reflection and boundary conditioning]
\label{prop:reflection-stability}
Let \(I\Subset\R\).
\begin{enumerate}[label=\textnormal{(\roman*)},leftmargin=1.5em]
\item \(\Theta\) maps \(\mcA(I)\) onto \(\mcA(-I)\).
\item Conditional expectation onto \(\mcF_0\) maps
\(\mcA_{\mathrm{loc}}\cap L^2(\mathbb P)\) into \(L^2(\Omega,\mcF_0,\mathbb P)\).
\end{enumerate}
\end{proposition}

\begin{proof}
Item \textnormal{(i)} was already recorded in
Proposition~\ref{prop:local-algebra}\textnormal{(iii)}.  For
item \textnormal{(ii)}, conditional expectation is the orthogonal projection
from \(L^2(\mathbb P)\) onto the closed subspace
\(L^2(\Omega,\mcF_0,\mathbb P)\), hence preserves \(L^2\) and measurability
with respect to \(\mcF_0\).
\end{proof}

\begin{lemma}[Future-functional Markov property on the stationary path space]
\label{lem:future-markov}
Let \(G^+\in L^\infty(\Omega,\mcF_+,\mathbb P)\), and set
\[
g(y):=\E_{\mathbb P}[G^+\mid Y_0=y],
\]
so that \(g\in L^\infty(\pi_\infty)\) and
\[
\E_{\mathbb P}[G^+\mid\mcF_0]=g\circ Y_0.
\]
Then for every \(t>0\),
\[
\E_{\mathbb P}[\tau_t G^+\mid\mcF_0]
=
(T_t g)\circ Y_0.
\]
Equivalently, under the identification
\(L^2(\pi_\infty)\cong L^2(\Omega,\mcF_0,\mathbb P)\) through
\(h\mapsto h\circ Y_0\),
\[
\E_{\mathbb P}[\tau_t G^+\mid\mcF_0]
=
T_t\!\bigl(\E_{\mathbb P}[G^+\mid\mcF_0]\bigr).
\]
\end{lemma}

\begin{proof}
First consider a bounded cylinder future observable of the form
\[
G^+(\omega)=h\bigl(Y_{s_1}(\omega),\dots,Y_{s_k}(\omega)\bigr),
\qquad
0<s_1<\cdots<s_k,
\]
with \(h:(\R^8)^k\to\C\) bounded and Borel.  Then
\[
\tau_t G^+
=
h\bigl(Y_{t+s_1},\dots,Y_{t+s_k}\bigr).
\]
By repeated use of the Markov property of the stationary path law imported in
Proposition~\ref{prop:imports}\textnormal{(a)--(b)}, there exists a bounded
Borel function \(g_h:(\R^8)\to\C\) such that
\[
g_h(y)
:=
\E_y\!\left[h\bigl(Y_{s_1},\dots,Y_{s_k}\bigr)\right]
\]
and
\[
\E_{\mathbb P}[G^+\mid\mcF_0]=g_h(Y_0).
\]
Applying the same Markov property at time \(t\) gives
\[
\E_{\mathbb P}[\tau_t G^+\mid\mcF_0]
=
\E_{Y_0}\!\left[h\bigl(Y_{t+s_1},\dots,Y_{t+s_k}\bigr)\right]
=
T_t g_h(Y_0).
\]
Hence the stated identity holds for all bounded future cylinder observables.

Let \(\mathcal M\) be the collection of all bounded \(\mcF_+\)-measurable
observables \(G^+\) for which the identity holds.  The cylinder future
observables form an algebra, are closed under complex conjugation, and generate
\(\mcF_+\) because \(\Omega=C(\R;\R^8)\) and \(\mcF_+=\sigma(Y_s:s>0)\) is
generated by the coordinate maps at positive rational times.  Moreover,
\(\mathcal M\) is a monotone class: if \(0\le G_n^+\uparrow G^+\) with each
\(G_n^+\in\mathcal M\), then bounded convergence for conditional expectation yields
\[
\E_{\mathbb P}[\tau_t G_n^+\mid\mcF_0]\uparrow
\E_{\mathbb P}[\tau_t G^+\mid\mcF_0]
\]
almost surely, and likewise
\[
\E_{\mathbb P}[G_n^+\mid\mcF_0]\uparrow
\E_{\mathbb P}[G^+\mid\mcF_0].
\]
Since \(T_t\) is positivity preserving and bounded on
\(L^\infty(\pi_\infty)\cap L^2(\pi_\infty)\), one also has
\[
T_t\!\bigl(\E_{\mathbb P}[G_n^+\mid\mcF_0]\bigr)\uparrow
T_t\!\bigl(\E_{\mathbb P}[G^+\mid\mcF_0]\bigr)
\]
almost surely after identification with time-zero observables.  Therefore
\(\mathcal M\) contains the bounded monotone class generated by the cylinder algebra.
By the monotone class theorem, \(\mathcal M\) contains every bounded
\(\mcF_+\)-measurable observable, proving the lemma.
\end{proof}

\begin{proposition}[Past--future boundary transfer identity]
\label{prop:transfer}
Let \(F^-\in L^\infty(\Omega,\mcF_-,\mathbb P)\cap L^2(\mathbb P)\) and
\(G^+\in L^\infty(\Omega,\mcF_+,\mathbb P)\cap L^2(\mathbb P)\).  Define the
centered boundary data
\[
f^-:=\E_{\mathbb P}[F^-\mid \mcF_0]-\E_{\mathbb P}[F^-],
\qquad
g^+:=\E_{\mathbb P}[G^+\mid \mcF_0]-\E_{\mathbb P}[G^+].
\]
Then for every \(t>0\),
\[
\E_{\mathbb P}\!\bigl[F^-\,(\tau_t G^+)\bigr]
-\E_{\mathbb P}[F^-]\E_{\mathbb P}[G^+]
=
\left\langle f^-,T_t g^+\right\rangle_{L^2(\pi_\infty)},
\]
where \(T_t\) is the stationary mean-field Markov semigroup acting on
\(L^2(\pi_\infty)\cong L^2(\Omega,\mcF_0,\mathbb P)\) through the time-zero
coordinate map \(h\mapsto h\circ Y_0\).
\end{proposition}

\begin{proof}
By Proposition~\ref{prop:hyperplane-factorization},
\[
\E_{\mathbb P}\!\bigl[F^-\,(\tau_t G^+)\bigr]
=
\E_{\mathbb P}\!\left[
\E_{\mathbb P}[F^-\mid \mcF_0]\,
\E_{\mathbb P}[\tau_t G^+\mid \mcF_0]
\right].
\]
By Lemma~\ref{lem:future-markov},
\[
\E_{\mathbb P}[\tau_t G^+\mid \mcF_0]
=
T_t\!\bigl(\E_{\mathbb P}[G^+\mid \mcF_0]\bigr).
\]
Subtracting the product of expectations yields
\[
\E_{\mathbb P}\!\bigl[F^-\,(\tau_t G^+)\bigr]
-\E_{\mathbb P}[F^-]\E_{\mathbb P}[G^+]
=
\left\langle
\E_{\mathbb P}[F^-\mid\mcF_0]-\E_{\mathbb P}[F^-],\,
T_t\!\bigl(\E_{\mathbb P}[G^+\mid\mcF_0]-\E_{\mathbb P}[G^+]\bigr)
\right\rangle_{L^2(\pi_\infty)},
\]
which is exactly the stated identity.
\end{proof}

\begin{corollary}[Reflected positive-time specialization]
\label{cor:reflected-transfer}
Let \(F,G\in L^\infty(\Omega,\mcF_+,\mathbb P)\cap L^2(\mathbb P)\), and set
\[
f_\Theta:=\E_{\mathbb P}[\Theta F\mid\mcF_0]-\E_{\mathbb P}[\Theta F],
\qquad
g:=\E_{\mathbb P}[G\mid\mcF_0]-\E_{\mathbb P}[G].
\]
Then for every \(t>0\),
\[
\E_{\mathbb P}\!\bigl[(\Theta F)\,(\tau_t G)\bigr]
-\E_{\mathbb P}[\Theta F]\E_{\mathbb P}[G]
=
\langle f_\Theta,T_t g\rangle_{L^2(\pi_\infty)}.
\]
\end{corollary}

\begin{proof}
If \(F\) is \(\mcF_+\)-measurable, then \(\Theta F\) is
\(\mcF_-\)-measurable by Proposition~\ref{prop:local-algebra}\textnormal{(iii)}.
Apply Proposition~\ref{prop:transfer} with \(F^-:=\Theta F\) and \(G^+:=G\).
\end{proof}

\begin{corollary}[Reduction of Euclidean correlation estimates to boundary semigroup estimates]
\label{cor:reduction}
For every
\[
F^-\in L^\infty(\Omega,\mcF_-,\mathbb P)\cap L^2(\mathbb P),
\qquad
G^+\in L^\infty(\Omega,\mcF_+,\mathbb P)\cap L^2(\mathbb P),
\]
any estimate on \(T_t\) acting on centered boundary observables transfers
directly to the corresponding mixed past--future Euclidean correlation estimate
by Proposition~\ref{prop:transfer}.  The reflected positive-time case is the
specialization recorded in Corollary~\ref{cor:reflected-transfer}.
\end{corollary}

\begin{proof}
This is immediate from Proposition~\ref{prop:transfer}.  For example, any
bound of the form
\[
\|T_t h\|_{L^2(\pi_\infty)}\le a_t\|h\|_{L^2(\pi_\infty)}
\]
implies
\[
\left|
\E_{\mathbb P}\!\bigl[F^-\,(\tau_t G^+)\bigr]
-\E_{\mathbb P}[F^-]\E_{\mathbb P}[G^+]
\right|
\le
a_t\,\|f^-\|_{L^2(\pi_\infty)}\,\|g^+\|_{L^2(\pi_\infty)}.
\]
\end{proof}

\section{Conclusion}

This section supplies the Euclidean ingredients used in the Osterwalder--Schrader verification.  It fixes
the Euclidean thermodynamic theory and proves the realization-specific
structural statements used in the subsequent verification:
\begin{enumerate}[leftmargin=1.5em]
\item the exact stationary Euclidean probability space;
\item the local bounded observable algebras on time windows;
\item the continuum-indexed Schwinger functionals;
\item the thermodynamic-first identification of the gauge sector;
\item the hyperplane factorization, reflection stability, the exact past--future
transfer identity, and its reflected specialization.
\end{enumerate}
The remaining work is axiomatic: verifying OS0--OS4,
reconstructing the Minkowski theory, and converting the constructive spectral
gap into the mass gap.

\begin{thebibliography}{9}
\bibitem{convergence}
\emph{01\_convergence}, companion paper in the series.

\bibitem{continuumqft}
\emph{03b\_continuum\_yang\_mills}, companion paper in the series.

\bibitem{parthasarathy1967}
K.~R. Parthasarathy,
\emph{Probability Measures on Metric Spaces},
Academic Press, 1967.
\end{thebibliography}

\end{document}
